\section{Step 24 Advance the outer input by one sample}\label{fir-s24}

Advance \texttt{a0} and \texttt{a2} by two, consume one output and repeat \texttt{scalar\_window} while count remains. The new input window overlaps seven samples of the previous window. Its private walkers and accumulator restart, while \texttt{v1} keeps the call’s accumulated clipping count. The scalar path now implements every supported count, including asymmetric impulse responses and maximum sums.

\Needspace{12\baselineskip}
\subsection{Write and inspect}
Replace the temporary branch and nop in \texttt{next\_window}. Keep surrounding labels and region markers.

\Needspace{8\baselineskip}
\MARSCodeLabel{MIPS assembly}
\begin{minted}[samepage=true]{mars-idea}
    daddiu $a0, $a0, 2
    daddiu $a2, $a2, 2
    addiu $a3, $a3, -1
    bne $a3, $zero, scalar_window
    nop
\end{minted}
Use recipe \texttt{s24-1}, \texttt{s24-2}, \texttt{s24-3} from the companion construction inputs. Pause at \texttt{filter\_ready} before comparing the partial state.

\subsection{Expected observation}
The default sixteen-output impulse returns eight samples of 2000 then zeros. The asymmetric two-output fixture yields 4096 then 3584. Sixty-four maximum positive sums return sixty-four clipping events.

\Needspace{11\baselineskip}
\subsection{Check your understanding}
Why advance the input by two rather than sixteen bytes?

\begin{itemize}
\item \textbf{A} Each next output slides the eight-tap window by one sample.
\item \textbf{B} The window has only one input byte.
\item \textbf{C} The coefficients advance automatically with \texttt{a0}.
\end{itemize}
